\begin{helper}
Let \(U\) be a finite union of witness events \(W_i\), indexed by a finite
linearly ordered type:
\[
  U=\bigcup_i W_i .
\]
Then there is a canonical first-witness refinement \(P_i\) of this union,
defined by
\[
  P_i=\{a:a\in W_i\hbox{ and }a\notin W_j\hbox{ for every }j<i\}.
\]
The refined pieces still cover \(U\),
\[
  U=\bigcup_i P_i,
\]
and the family \((P_i)_i\) is pairwise disjoint.
\end{helper}
\begin{proof}
Define \(P_i\) by the displayed formula.  First we prove that the \(P_i\)
cover \(U\).  If \(a\in P_i\) for some \(i\), then \(a\in W_i\) by the first
clause in the definition of \(P_i\), and hence \(a\in U\) because
\(U=\bigcup_i W_i\).

Conversely suppose \(a\in U\).  Since \(U=\bigcup_i W_i\), the finite set
\[
  S(a)=\{i:a\in W_i\}
\]
is nonempty.  Choose its least element \(i_0\) in the fixed finite linear
order.  Then \(a\in W_{i_0}\).  If \(j<i_0\), minimality of \(i_0\) in
\(S(a)\) gives \(j\notin S(a)\), hence \(a\notin W_j\).  Therefore
\(a\in P_{i_0}\), so \(a\) lies in the union of the refined pieces.

Now take two distinct indices \(i\ne j\).  Since the order is linear, either
\(i<j\) or \(j<i\).  If \(a\in P_i\cap P_j\) and \(i<j\), then the
first-witness exclusion in the definition of \(P_j\) says \(a\notin W_i\),
whereas \(a\in P_i\) says \(a\in W_i\), a contradiction.  The case \(j<i\)
is symmetric.  Thus \(P_i\cap P_j=\emptyset\) whenever \(i\ne j\), which is
pairwise disjointness.

Finally, intersecting every event with a fixed cell preserves both the union
identity and disjointness: \((\bigcup_i P_i)\cap C=\bigcup_i(P_i\cap C)\),
and subsets of disjoint sets remain disjoint.  Pairwise disjointness implies
pairwise almost-disjointness for any measure, since the intersection of two
distinct refined pieces is empty and therefore has measure zero.
\end{proof}
